%
\againframe{plan}
%
\begin{frame}{Lecture Plan}
	\begin{learningobjectives}
		Learn the concept of 
		\begin{itemize}
			\item electrical conductivity,
			\item mobility, and
			\item relaxation time
		\end{itemize}
	\end{learningobjectives}
\end{frame}
%
\begin{frame}{Nature of Metal}
	%
	\begin{figure}
		\includegraphics[width=\textwidth]{metals_nature}
	\end{figure}
	%
\begin{columns}
	\begin{column}[t]{0.6\textwidth}
		\begin{enumerate}
			\item Lustre (Shine)
			\item Solid with high \SI{~1000}{\celsius} melting points
			\item Malleable (capable of being shaped)
			\item Good electrical conductor
			\item Good thermal conductor
			\item Ductile (easy to draw wires)
		\end{enumerate}
	\end{column}
	%
	\begin{column}[t]{0.38\textwidth}
		\begin{block}{Chemistry}
			\begin{itemize}
				\item Metallic bonding
				\item Screening
			\end{itemize}
			%
			\begin{figure}
				\centering
				\includegraphics[width=0.4\textwidth]{metals_screening}
			\end{figure}
		\end{block}
	\end{column}
\end{columns}
	% 

\end{frame}
%
\begin{frame}{Metallic bonding \(\leftrightarrow\) electron gas}
	%
		\begin{figure}
			\centering
			\includegraphics[width=0.9\textwidth]{metals_free_electron_gas}
			\caption{Valence electrons to electron ``gas''}
		\end{figure}
	%
		\begin{keyinsight}
		The properties of electron ``gas'' determines the nature of metal.
	\end{keyinsight}
	%
\end{frame}
%
\begin{frame}{Electron theories of metals}
\begin{columns}
	\begin{column}{0.3\textwidth}
		\begin{enumerate}
			\item Classical free electron theory
			\item Quantum free electron theory \hfill [M2 U3]
			\item Quantum band theory \hfill [M2 U3]
		\end{enumerate}
	%
	\end{column}
	%
	\begin{column}{0.68\textwidth}
			\begin{figure}
			\centering
			\includegraphics[width=\textwidth]{metals_discovery_timeline}
		\end{figure}
	\end{column}
\end{columns}
\end{frame}
%
